% Chapter 1 is independently editable. Shared facts are in config/facts.tex.
\chapter{Introduction}
\label{ch:introduction}

\ProjectTitle{} is a local web application for finding visually similar clothing
from an image. A visitor can search with a complete photograph or select the
garment of interest before searching. A local administrator can maintain a
separate product collection through the same application. This chapter establishes
the problem, objectives, scope, operating environment, and limits of the project.
The implementation scope described here corresponds to presentation release
\texttt{\ApplicationRelease}.

\section{Study Background}
\label{sec:study-background}

A clothing photograph can communicate a garment's appearance more directly than
a short product description. A person may recognise a jacket's silhouette or a
dress's pattern while being unsure which words a catalogue uses for it. Visual
retrieval addresses this situation by using an image to rank catalogue candidates.
Fashion retrieval research also distinguishes recognising a clothing category
from identifying the same item across different photographs. DeepFashion, for
example, provides clothing annotations and image correspondences across consumer,
street, and shop settings~\cite{Liu2016DeepFashion}. This distinction matters:
finding another dress is a different task from finding the same dress.

An image encoder maps a photograph to a numerical representation, or embedding,
which can be compared with stored catalogue representations. SigLIP2 is a family
of pretrained vision-language encoders developed for visual representation
learning and related tasks~\cite{Tschannen2025SigLIP2}. Our system uses its image
encoder to describe the query and catalogue images. We use pretrained model
weights; developing a new foundation model is outside the project.

A complete photograph may contain a shirt, trousers, shoes, a person, and
background objects. It does not necessarily tell the retrieval system which item
the user wants. Promptable segmentation provides a way to express that choice.
SAM2 supports interactive segmentation through user prompts, and its official
implementation supplies the improved SAM2.1 checkpoints used by this
project~\cite{Ravi2024SAM2,MetaSAM2Repository}. In our application, the user can
refine a selected region and inspect the resulting cutout before submitting it
for retrieval.

The final stage compares the query embedding with catalogue embeddings. FAISS
provides vector-search tools for this task~\cite{Johnson2017FAISS}. The application
uses \texttt{IndexFlatIP}, an exact inner-product index; normalising the vectors
allows this ranking to represent cosine similarity~\cite{FaissIndexes}.
Segmentation, encoding, indexing, and the browser interface must therefore work
together: a relevant representation of the wrong garment still answers the
wrong question.

\section{Problem Statement}
\label{sec:problem-statement}

The problem addressed by this project is how to retrieve relevant clothing from
a bounded catalogue when a query photograph may contain several competing items.
A full-image embedding represents the whole input, whereas the intended target
may occupy only part of it. Lighting, pose, background, occlusion, and differences
between query and catalogue photography can also change the representation.

The system must let the user communicate the target, prepare the query
consistently with catalogue images, rank the available candidates, and present
results that the user can inspect. It must also support full-image search when
selection is unnecessary. The selected cutout is therefore part of the search
workflow, rather than an invisible intermediate result. Category controls allow
the user to correct automatic routing or search across categories.

Catalogue coverage creates a separate limitation. If an item is absent, a
similarity search can return a related garment without returning the same
commercial product. The project consequently treats visual similarity,
category relevance, and exact product identity as separate concepts. A displayed
similarity score is not a calibrated probability that two photographs show the
same product. Evaluating identity requires verified product relationships and
different query and gallery photographs.

\section{Objectives}
\label{sec:objectives}

\subsection{General Aim}
The aim is to build and evaluate a local visual clothing search application that
combines interactive garment selection with embedding-based catalogue retrieval.

\subsection{Research Objectives}
\begin{enumerate}[leftmargin=1.8em,itemsep=0.35em]
  \item Compare pretrained visual encoders under documented gallery, query, and
    relevance definitions.
  \item Compare full-image and selected-garment retrieval, documenting the
    selection prompts and experimental conditions.
  \item Assess preprocessing and category-routing choices against a retained
    baseline before adopting a change.
  \item Report task-specific retrieval measures and evaluate product identity
    separately when verified alternate photographs are available.
\end{enumerate}

\subsection{Application Objectives}
\begin{enumerate}[leftmargin=1.8em,itemsep=0.35em]
  \item Provide image input, query preview, full-image search, SAM2 refinement,
    and manual selection within one browser workflow.
  \item Apply consistent image encoding, retrieve ranked candidates, and expose
    automatic and user-selected category controls.
  \item Support product creation, editing, and deletion with photographs,
    descriptive metadata, price, and availability.
  \item Protect local administration and catalogue recovery, and maintain a
    presentation release supported by source control and functional verification.
\end{enumerate}

\section{Scope of Study}
\label{sec:scope}

The application processes still images of clothing. Its principal workflow
combines SAM2.1 garment selection, SigLIP2 image embeddings with
\EmbeddingDimensions{} dimensions, and FAISS retrieval behind a FastAPI service.
Manual rectangle, polygon, and freehand selection are available when the user
prefers direct control. SAM2 interaction can be refined with positive and negative
points. The accepted query or cutout is previewed before search.

The original research gallery contains \TrainImageCount{} images in ten clothing
categories from the \texttt{clothing-dataset-small} collection~\cite{GrigorevClothingDataset}.
The dataset's named training partition serves as the retrieval gallery; it does
not imply that the project trained SigLIP2 or SAM2 on these images.
The \ValidationImageCount{} validation images and \TestImageCount{} test images are
kept outside that gallery for evaluation. These category-labelled partitions do
not establish that every image represents a unique commercial product.

A second, independent catalogue contains products added through the management
page. Each product can have one to five photographs, a name, an optional category
and price, a currency, and an availability state. The main search page can search
this collection using the same full-image and selection tools. Prices describe
the products and do not determine their visual-similarity ranking. Updating this
collection does not add its images to the original research gallery.

Visitors can search without an account. One local administrator controls product
management, catalogue export and restore, and account settings. Backups cover the
added-product collection and its associated images and embeddings; they are
separate from the research dataset, model checkpoints, and administrator account.
The application runs on the team's computer through a local browser connection.
Table~\ref{tab:scope} summarises the main boundaries.

\begin{table}[!htbp]
  \centering
  \caption{Implemented scope and deferred capabilities}
  \label{tab:scope}
  \begingroup\small\setstretch{1.1}\renewcommand{\arraystretch}{1.18}
  \begin{tabularx}{\linewidth}{@{}Y Y@{}}
    \toprule
    \textbf{Included in the current scope} & \textbf{Outside the current scope} \\
    \midrule
    Still-image search with full-image or selected-garment queries &
      Video retrieval and a native mobile application \\ \addlinespace
    Original research gallery and an independent added-product collection &
      An automatically collected city-wide or multi-store catalogue \\ \addlinespace
    One local administrator and visitor search &
      Seller accounts, store ownership, and permissions for multiple businesses \\ \addlinespace
    Product photographs, descriptive metadata, price, and availability &
      Orders, payment, checkout, and inventory transactions \\ \addlinespace
    Local backup, restore preview, and recovery copies &
      Managed cloud hosting and distributed storage \\ \addlinespace
    Documented category retrieval and limited reviewed identity studies &
      Guaranteed exact-SKU matching or general consumer-photo accuracy \\
    \bottomrule
  \end{tabularx}
  \endgroup
\end{table}

\section{Hardware and Software Requirements}
\label{sec:requirements}

\subsection{Verified Hardware Environment}
Table~\ref{tab:hardware} records the computer available for development and the
local presentation, checked on \EnvironmentCheckDate. It is a reference
environment in which the application has been exercised, not a measured minimum
specification. GPU acceleration is used for neural inference when CUDA is
available. The installed FAISS package performs vector search on the CPU.

\begin{table}[!htbp]
  \centering
  \caption{Verified local presentation environment}
  \label{tab:hardware}
  \begingroup\small\setstretch{1.1}\renewcommand{\arraystretch}{1.18}
  \begin{tabularx}{\linewidth}{@{}>{\raggedright\arraybackslash}p{0.24\linewidth}Y@{}}
    \toprule
    \textbf{Component} & \textbf{Verified configuration or requirement} \\
    \midrule
    Processor & \ReferenceProcessor \\ \addlinespace
    Memory & \ReferenceMemory \\ \addlinespace
    Graphics & \ReferenceGPU \\ \addlinespace
    Operating system & \ReferenceOS \\ \addlinespace
    Local storage & Space for model checkpoints, dataset images, vector indices,
      application databases, and retained backup files \\ \addlinespace
    User access & A local browser, keyboard, and pointing device; model and data
      files prepared before the presentation \\
    \bottomrule
  \end{tabularx}
  \endgroup
\end{table}

The implementation contains a CPU fallback, but its latency depends on the model,
input, and hardware. We therefore do not infer a performance guarantee for an
untested machine. Once the required assets are available locally, the presentation
launcher uses the local model files without downloading them during startup.

\FloatBarrier
\subsection{Software Environment}
The backend uses Python \PythonRuntimeVersion{} in an isolated virtual
environment. PyTorch and Transformers load the neural models, FAISS ranks vectors,
and FastAPI with Uvicorn serves the application. The frontend uses HTML, CSS, and
JavaScript, while SQLite stores the added-product collection and account data.
Table~\ref{tab:software} identifies the main responsibilities.

\begin{table}[!htbp]
  \centering
  \caption{Principal software components}
  \label{tab:software}
  \begingroup\small\setstretch{1.1}\renewcommand{\arraystretch}{1.18}
  \begin{tabularx}{\linewidth}{@{}>{\raggedright\arraybackslash}p{0.29\linewidth}Y@{}}
    \toprule
    \textbf{Component} & \textbf{Role in the project} \\
    \midrule
    Python and virtual environment & Backend execution and dependency isolation \\ \addlinespace
    PyTorch and Transformers & Local SigLIP2 and SAM2.1 inference \\ \addlinespace
    FAISS CPU and NumPy & Normalised vector storage, comparison, and ranking \\ \addlinespace
    FastAPI and Uvicorn & Image, segmentation, search, and administration services \\ \addlinespace
    SQLite & Product records, photographs, embeddings, and local account storage \\ \addlinespace
    HTML, CSS, and JavaScript & Browser search, selection, and administration interface \\ \addlinespace
    Git and LaTeX & Version control and editable academic documentation \\
    \bottomrule
  \end{tabularx}
  \endgroup
\end{table}

The repository records application dependencies, while checkpoints, the dataset,
local databases, and backups are separate runtime resources. Access to source
code alone is therefore insufficient to reproduce the complete local environment;
the setup instructions must also identify these resources.

\FloatBarrier
\section{Significance of Project}
\label{sec:significance}

The practical value of the project is the connection between a user's visual
intent and an inspectable retrieval process. Selection lets the user indicate
which garment matters. Previewing the submitted cutout makes selection errors
visible before retrieval, and category correction provides another way to refine
a search without changing the underlying encoder.

Product management gives the search system a maintained collection to search.
A catalogue operator can update photographs and descriptive information through
the browser, while visitor search remains available without administrative access.
Backup and recovery support the continued use of that local collection. This
demonstrates a path towards store-specific visual search without claiming that
multi-store deployment has already been implemented.

The research value lies in evaluating the components under explicit conditions.
Encoder comparison, garment selection, preprocessing, and category routing can
be investigated separately. Reporting the data and relevance definition alongside
each result makes it possible to discuss both improvements and failures without
treating all measures as one notion of accuracy. The project contributes an
integrated, evaluated prototype rather than a newly trained foundation model.

\section{Constraints and Critical Assumptions}
\label{sec:constraints}

\begin{description}[leftmargin=0pt,labelindent=0pt,font=\normalfont\bfseries,itemsep=0.6em]
  \item[Catalogue coverage.]
    Retrieval is limited to available catalogue content. A relevant-looking
    result does not establish that the requested commercial product is present.
  \item[Product identity.]
    Original catalogue identifiers are derived from image files. Category labels
    and file identifiers are insufficient to establish exact-SKU ground truth.
    Reviewed alternate views support only the scope of their recorded study.
  \item[Selection quality.]
    A mask can include background or omit part of a garment. We assume that the
    user inspects and, when necessary, refines the selection or uses a manual tool.
    Successful retrieval does not establish pixel-level mask accuracy.
  \item[Pretrained representations.]
    The encoders may emphasise category, shape, texture, and colour differently.
    Visual similarity is not calibrated identity confidence, and a better
    category match need not be the closest commercial product.
  \item[Evaluation comparability.]
    Results from different queries, gallery exclusions, preprocessing versions,
    or relevance definitions must remain separate. Functional test counts measure
    exercised software behaviour, not retrieval accuracy.
  \item[Local operation and data.]
    The current deployment assumes the team's computer, available local models,
    and one administrator. The database and its backups require local protection;
    the implemented account controls do not establish readiness for public,
    multi-tenant hosting.
\end{description}

The original catalogue index remains the active presentation baseline. Catalogue
cropping and other ranking changes are experimental choices, not assumed
improvements. Any later change must be compared with the retained baseline
under an appropriate evaluation protocol.

\section{Project Management}
\label{sec:project-management}

Development is organised into problem definition, data preparation, model
evaluation, implementation, integration, verification, and documentation.
Experimental comparisons iterate within these stages. A proposed improvement is
checked against its baseline before promotion to the presentation version.
This approach accommodates both software defects and research findings that do
not support adopting a change.

Git records the application source and presentation releases. The current
documentation uses \texttt{\ApplicationRelease} as its implementation reference.
Table~\ref{tab:milestones} summarises the work packages at the time of this
chapter's preparation. Completed implementation and recorded experiments are
distinguished from the remaining submission work.

\begin{table}[!htbp]
  \centering
  \caption{Project work packages and current status}
  \label{tab:milestones}
  \begingroup\small\setstretch{1.1}\renewcommand{\arraystretch}{1.18}
  \begin{tabularx}{\linewidth}{@{}>{\raggedright\arraybackslash}p{0.08\linewidth}Y >{\raggedright\arraybackslash}p{0.18\linewidth}@{}}
    \toprule
    \textbf{ID} & \textbf{Work package and output} & \textbf{Status} \\
    \midrule
    M1 & Dataset preparation and baseline encoder comparison & Recorded \\ \addlinespace
    M2 & Full-image search, SAM2 refinement, and manual selection & Implemented \\ \addlinespace
    M3 & Shared preprocessing, category controls, and query preview & Implemented \\ \addlinespace
    M4 & Product editing, local administrator access, backup, and restore & Implemented \\ \addlinespace
    M5 & Functional verification and reviewed identity pilot & Recorded \\ \addlinespace
    M6 & LaTeX report and chapter review & In progress \\ \addlinespace
    M7 & Final presentation, poster, brochure, video, and submission review & Pending \\
    \bottomrule
  \end{tabularx}
  \endgroup
\end{table}

\FloatBarrier
The remaining academic materials will be derived from the same verified scope
and evidence. The university's announced discussion period is 4--7 November
2026. Document preparation proceeds chapter by chapter so that terminology,
references, diagrams, and results can be reviewed before the report is submitted
for supervisor and faculty approval. Later feedback may require revisions to the
report or a separately tested application change.

This chapter has established the intended users, retrieval problem, implemented
scope, operating environment, and principal limits. The next chapter reviews
the literature supporting the encoder, segmentation, and retrieval choices.
